\documentclass{article}
\usepackage[UTF8]{ctex}
\usepackage{amsmath,amssymb,amsfonts}
\usepackage{enumitem}
\usepackage{geometry}
\geometry{a4paper, margin=1in}
% === 错题本排版优化（不参与题目标识） ===
\setlength{\parindent}{0pt}
\setlength{\parskip}{0.35em plus 0.10em minus 0.08em}
\setlength{\emergencystretch}{3em}
\setlist{itemsep=0.25em, topsep=0.40em, parsep=0pt}
\allowdisplaybreaks[2]
% === 错题本排版优化结束 ===
\begin{document}
\section*{中国矿业大学2024年数学分析真题}
\section*{22．中国矿业大学2024年数学分析试卷}

\subsection{题目1}
1．（15 分）设定义在 \([a, b]\) 上的连续函数 \(f(x)\) 满足 \(f(a)=f(b)\) ，且 \(f_{+}^{\prime}(a) f_{-}^{\prime}(b)>0\) 。证明：在 \((a, b)\)内至少有一点 \(\xi\) ，使得 \(f(\xi)=f(a)\) ．

\subsection{题目2}
2．（15 分）叙述并证明 \(\lim _{x \rightarrow+\infty} f(x)=+\infty\) 的归结原则。

\subsection{题目3}
3．（15 分）叙述有限覆盖定理和致密性定理，并用有限覆盖定理证明致密性定理．

\subsection{题目4}
4．（15 分）求极限 \(\lim _{x \rightarrow 0} \frac{1}{x}\left(\frac{1}{x}-\cot x\right)\) ．

\subsection{题目5}
5．（15 分）设 \(f(x)\) 与 \(g(x)\) 均是定义在 \([a, b]\) 上的可积函数，证明：函数 \(f(x) g(x)\) 也在 \([a, b]\) 上可积．

\subsection{题目6}
6．（15 分）求 \(I=\int_{0}^{1} t g(t) \mathrm{d} t\) ，其中 \(g(t)=\int_{t^{2}}^{1} e^{-x^{2}} \mathrm{~d} x\) ．

\subsection{题目7}
7．（15 分）求符号函数 \(f(x)=\operatorname{sgn} x\) 在 \((-\pi, \pi)\) 上的傅里叶级数并判断其敛散性．

\subsection{题目8}
8．（15 分）证明当 \(\alpha>0\) 时，函数项级数 \(\sum_{n=1}^{\infty} \frac{\sin n x}{n^{2+\alpha}}\) 在 \((-\infty,+\infty)\) 上一致收敛，且其和函数 \(S(x)\) 在\((-\infty,+\infty)\) 内具有连续的导函数．

\subsection{题目9}
9．（15 分）设二元函数组 \(\left\{\begin{array}{l}u=f(x, y) ; \\ v=g(x, y) .\end{array}\right.\) 将 \(x O y\) 平面上的点集 \(E\) 映射成 \(u O v\) 平面上的点集D．\(f(x, y)\) 和 \(g(x, y)\) 为点集 \(E\) 上的一致连续函数，且 \(h(u, v)\) 为点集 \(D\) 上的一致连续函数，证明：复合函数 \(h(f(x, y), g(x, y))\) 为点集 \(E\) 上的一致连续函数．

\subsection{题目10}
10．（15 分）设 \(I(\alpha)=\int_{0}^{+\infty} e^{-\alpha x} \frac{\sin \beta x}{x} \mathrm{~d} x\) ，其中 \(\beta>0\) 为常数．证明：

\begin{itemize}
  \item[（1）] （7 分）\(I(\alpha)\) 在 \((0,+\infty)\) 上连续．
  \item[（2）] （5 分）\(I(\alpha)\) 在 \((0,+\infty)\) 上可导并求 \(I(\alpha)\) 的表达式．
  \item[（3）] （3）分）利用 \(I(\alpha)\) 的表达式求 \(D(\beta)=\int_{0}^{+\infty} \frac{\sin \beta x}{x} \mathrm{~d} x\) 的值．
\end{itemize}
\end{document}
